\section{Product fillings and their limits}\label{sec:product-fillings}

\subsection{The cycles created by a filling}

A meridian of a solid torus bounds a disk.  A longitude runs once around
its core.  Gluing two solid tori therefore kills two boundary cycles,
but it also joins their three-dimensional interiors.  The first operation
controls torsion.  The second supplies a top-dimensional cycle.
Both operations must be retained in an integral chain calculation.

We use singular homology with integer coefficients unless another ring
is displayed.  An integral $j$-chain is a finite integer sum of oriented
singular $j$-simplices.  Its boundary is the alternating sum of its faces.
A cycle has zero boundary, and two cycles represent the same homology
class when their difference is a boundary.  Thus, for a chain complex
$(C_*,\partial)$,
\[
 H_j(C_*)=\ker(\partial:C_j\to C_{j-1})/
                 \operatorname{im}(\partial:C_{j+1}\to C_j).
\]

A chain map commutes with boundaries.  A chain homotopy between maps
$f,g$ is a collection $h_j:C_j\to D_{j+1}$ with
$f-g=\partial h+h\partial$.  Homotopic chain maps induce equal maps
on homology.

Write $\T^k=(\mathbb R/\Z)^k$.  Give $T=\T^2$ an ordered basis
$(x,c)$ of oriented coordinate circles and orient it by $x\cdot c=1$.
Fix integers
\begin{equation}\label{prod:eq:data}
 r,s>0,\qquad m,n,b_0\in\Z,\qquad
 \nu=n-sb_0,\qquad \Delta=r\nu+sm.
\end{equation}
The two proposed meridians are
\begin{equation}\label{prod:eq:slopes}
                  \mu_0=rx-mc,\qquad \mu_1=sx+\nu c.
\end{equation}

A lattice vector is \emph{primitive} when its coordinates have greatest
common divisor one.  This means that it generates a direct summand:
integers $a,b$ with $au+bv=1$ complete $(u,v)$ to a unimodular basis.
A meridian must be primitive, since the boundary homomorphism
$H_1(\partial(D^2\times S^1))\to H_1(D^2\times S^1)$ has
kernel generated by the disk boundary.

First assume
\begin{equation}\label{prod:eq:effective}
                         \gcd(r,m)=\gcd(s,n)=1.
\end{equation}
Take $V_i\cong D^2\times S^1$ with meridian $\mu_i$.
Use linear torus markings on both boundaries.
Identify $T$ with $\partial V_0$ with its boundary orientation, and glue
$V_1$ by an orientation-reversing linear boundary identification.
The result is an oriented closed manifold $M$.

Within these linear markings, the two meridians determine the construction
up to homeomorphism.
Indeed, the change between two choices of longitude preserves the meridian
and extends across the disk by rotating it along the core circle.

The boundary classes $x_0=x$ and $x_1=-x+b_0c$ satisfy
$x_0+x_1=b_0c$.  The meridian relations become
$rx_0=mc$ and $sx_1=nc$.  These equations explain the section integer
$b_0$ in \eqref{prod:eq:data}.  They also locate the regular fiber $c$ in the
same integral boundary lattice as the filling slopes.

\begin{example}\label{prod:ex:eight}
Let $(r,s;m,n;b_0)=(4,4;1,1;0)$.  The meridians are
$4x-c$ and $4x+c$.  Subtracting their relations gives $2c=0$,
while the first gives $c=4x$.  Therefore $8x=0$.
The resulting first homology is cyclic of order eight.
We shall compute $H_2(M)=0$ and $H_3(M)=\Z$ from the same gluing,
then exhibit three independent order-eight classes in each of degrees
two and three of $M\times\T^3$.
\end{example}

\subsection{A cellular gluing complex}

A circle has one vertex and one oriented edge with zero boundary.
The product cell structure on $T$ has one vertex, edges $x,c$, and
one face $x\times c$.  Opposite edges of its square cancel, so its
cellular boundaries are zero.
The corresponding chain groups are $\Z$, $\Z^2$, and $\Z$ in
degrees zero, one, and two.

We recall why cellular chains compute integral homology in the finite
cell spaces used here.  Filter such a space by its skeleta $X^{(j)}$.
Collapsing $X^{(j-1)}$ in $X^{(j)}$ gives a wedge of oriented
$j$-spheres, one for each $j$-cell.  The relative homology is therefore
free on those cells in degree $j$, and zero in other degrees.

The connecting maps for consecutive skeleta give the cellular boundary.
Exactness for each triple of consecutive skeleta gives $\partial^2=0$.
The same exact sequences, inductively over the finite skeleta, identify
its cycles modulo boundaries with singular homology.

Relative homology here means the homology of the quotient of singular
chains by the chains in the specified subspace.  These statements also
show that collapsing a contractible edge joining distinct vertices
preserves homology.

Orient the core circles so that inclusion and retraction give the winding
homomorphisms
\begin{equation}\label{prod:eq:core-maps}
 f_0(x)=m,\quad f_0(c)=r,\qquad
 f_1(x)=-\nu,\quad f_1(c)=s.
\end{equation}
Their kernels are respectively $\Z\mu_0$ and $\Z\mu_1$.
Each homomorphism is onto by \eqref{prod:eq:effective}.
Concretely, in coordinates $(u,v)\in T$, the two maps to $\mathbb R/\Z$
are $mu+rv$ and $-\nu u+sv$.

The \emph{double mapping cylinder} of $f_0,f_1$ is
\begin{equation}\label{prod:eq:pushout}
 P=\bigl(S^1_0\amalg(T\times[0,1])\amalg S^1_1\bigr)/
       ((z,0)\sim f_0(z),\ (z,1)\sim f_1(z)).
\end{equation}
It keeps a cylinder between the two maps instead of identifying their
values directly.  This retention of the boundary cylinder is what
preserves the higher cells.

\begin{lemma}\label{prod:lem:geometric-cylinder}
Under \eqref{prod:eq:effective}, the space $P$ is homotopy equivalent to $M$.
\end{lemma}

\begin{proof}
A primitive winding map $T\to S^1$ becomes projection onto the second
circle after a unimodular change of boundary coordinates.
The mapping cylinder of $S^1\times S^1\to S^1$ is
$\operatorname{Cone}(S^1)\times S^1\cong D^2\times S^1$.
Its uncollapsed boundary is $T$, and its meridian is the kernel of the
winding map.  Divide the cylinder in \eqref{prod:eq:pushout} at its midpoint.
Its two halves are precisely these solid tori, with meridians
\eqref{prod:eq:slopes}.  This proves the assertion for the winding maps.

For cellular calculations, choose cellular maps homotopic to these
winding maps.  Such a choice leaves \eqref{prod:eq:pushout} unchanged up to
homotopy equivalence.  To see this directly, attach the homotopy on a short
collar at each endpoint of $T\times[0,1]$.  Reversing the homotopy gives
a map in the other direction.  Their composites are homotopic to the
identity by contracting the consecutive forward and reverse homotopies.
The middle part of the cylinder stays fixed.
\end{proof}

The cellular maps can be constructed without further topological input.
Send the two edges around each target circle with the prescribed integer
windings.  The attaching word $xcx^{-1}c^{-1}$ of the torus face then
has total winding zero.  Its map to the circle lifts to a closed path
in $\mathbb R$, which contracts by linear interpolation.
This supplies the map on the face.
Two maps with equal edge windings are homotopic: their difference lifts
to $\mathbb R$ and contracts in the same way.

\begin{theorem}[Integral gluing chains]\label{prod:thm:chains}
The homology of $M$ is computed by the finite free chain complex
\begin{equation}\label{prod:eq:chain}
 0\longrightarrow\Z\mathrel{\mathop{\longrightarrow}^{0}}
 \Z^2\mathrel{\mathop{\longrightarrow}^{B}}\Z^2
 \mathrel{\mathop{\longrightarrow}^{0}}\Z\longrightarrow0,
 \qquad
 B=\begin{pmatrix}m&r\\ \nu&-s\end{pmatrix}.
\end{equation}
The displayed groups have degrees $3,2,1,0$.
The degree-two generators are the suspended boundary edges $x,c$.
The degree-one generators are the two oriented cores.
The degree-three generator comes from the suspended torus face.
\end{theorem}

\begin{proof}
The double mapping cylinder has the cells of the two target circles,
together with one cell of dimension $j+1$ for each $j$-cell of $T$.
Let $\sigma a$ denote a suspended cell.  Fix the suspension orientation by
\[
       \partial(\sigma a)=(f_0-f_1)(a)-\sigma(\partial a).
\]
This is the usual boundary of a prism with its orientation chosen by
the displayed formula.  The suspended vertex has boundary
$v_0-v_1$.  It forms a contractible edge joining the two vertices.

After collapsing this edge, one vertex and the two core edges remain.
The suspended edges have boundaries $(f_0-f_1)(x)$ and
$(f_0-f_1)(c)$, giving the columns of $B$.
The suspended face has zero boundary: the torus cellular boundary
is zero, and neither target circle has a two-cell.
Thus the cellular complex after the edge collapse is exactly
\eqref{prod:eq:chain}.  Lemma~\ref{prod:lem:geometric-cylinder} identifies its
homology with that of $M$.
\end{proof}

In particular, the determinant does not occur as a map out of degree
three.  The degree-three differential vanishes for the geometric reason
that the boundary face has nowhere to map in either core circle.
The top class can be oriented to agree with the fundamental cycle of $M$.

\subsection{Degree two, degree three, and prime coefficients}

\begin{theorem}\label{prod:thm:homology3}
For the effective filling,
\[
 H_0(M)=H_3(M)=\Z,\qquad
 (H_1(M),H_2(M))=
 \begin{cases}
 (\Z/|\Delta|\Z,0),&\Delta\ne0,\\
 (\Z,\Z),&\Delta=0.
 \end{cases}
\]
All groups in higher degrees vanish.
If $\Delta=0$, the primitive degree-two cycle in
\eqref{prod:eq:chain} is $r\,\sigma x-m\,\sigma c$, the suspension of $\mu_0$.
\end{theorem}

\begin{proof}
Choose $\alpha,\beta\in\Z$ with $m\alpha+r\beta=1$ and set
$q=s\beta-\nu\alpha$.  The unimodular matrix
\[
 Q=\begin{pmatrix}r&\alpha\\-m&\beta\end{pmatrix}
\]
has columns $\mu_0,\lambda_0$, where $\lambda_0=\alpha x+\beta c$.
Direct multiplication gives
\begin{equation}\label{prod:eq:reduction}
 BQ=\begin{pmatrix}0&1\\\Delta&-q\end{pmatrix},\qquad
 \begin{pmatrix}1&0\\q&1\end{pmatrix}
 BQ\begin{pmatrix}0&1\\1&0\end{pmatrix}
        =\begin{pmatrix}1&0\\0&\Delta\end{pmatrix}.
\end{equation}

The first diagonal summand $\Z\xrightarrow{1}\Z$ is contractible.
Deleting it is a chain homotopy equivalence: its contracting homotopy
sends its degree-one generator back to its degree-two generator.
The remaining complex has one generator in each degree from zero to
three and the single boundary $\partial e_2=\Delta e_1$.
This proves the group formulas.  At $\Delta=0$, the kernel of $B$
is generated by $\mu_0$, which is primitive by effectivity.
\end{proof}

If $\Delta=0$, reducing $r\nu+sm=0$ modulo $r$ gives $r\mid s$.
Reducing modulo $s$ gives $s\mid r$.  Hence $r=s$ and $\nu=-m$:
the two meridians coincide.  Geometrically, their disks join along this
common boundary to form a two-sphere.  This realizes the degree-two
class in Theorem~\ref{prod:thm:homology3}.
It also explains why determinant zero is a different topological case.

Cohomology reverses the arrows by applying $\operatorname{Hom}(-,R)$
to the free cellular chains, where $R$ is the coefficient ring.
In particular, for $\Delta\ne0$,
\[
 H^1(M;\Z)=0,\qquad H^2(M;\Z)=\Z/|\Delta|\Z,
 \qquad H^3(M;\Z)=\Z.
\]
The vanishing of integral $H_2$ therefore says nothing about the
vanishing of integral $H^2$.

\begin{proposition}[The exceptional primes]\label{prod:prop:prime}
Let $p$ be prime.  For $\Delta\ne0$,
\[
 H_1(M;\F_p)=H_2(M;\F_p)=
 \begin{cases}0,&p\nmid\Delta,\\\F_p,&p\mid\Delta.
 \end{cases}
\]
In all cases $H_0(M;\F_p)=H_3(M;\F_p)=\F_p$.
When $p\mid\Delta$, the connecting homomorphism for
$0\to\Z\xrightarrow{p}\Z\to\F_p\to0$ sends
\begin{equation}\label{prod:eq:bockstein}
 [e_2]\longmapsto (\Delta/p)[e_1]\in H_1(M;\Z).
\end{equation}
Its image has exact order $p$.
\end{proposition}

\begin{proof}
The unimodular transformations in \eqref{prod:eq:reduction} remain
invertible over $\F_p$.  Thus the only remaining differential is
multiplication by $\Delta$ in $\F_p$, proving the homology statements.

For the connecting map, lift the mod-$p$ cycle $e_2$ to its integral
cell.  Its boundary is $\Delta e_1$.  Divide this boundary by $p$
and take its integral homology class, giving \eqref{prod:eq:bockstein}.
Changing a lift by $p$ times a chain changes this class by a boundary.
The least positive multiple of $\Delta/p$ divisible by $\Delta$
is $p$, proving the last assertion, including negative $\Delta$.
\end{proof}

In Example~\ref{prod:ex:eight}, integral $H_2$ vanishes but
$H_2(M;\F_2)=\F_2$.  Its connecting map has image generated by
$4[e_1]$ in $\Z/8\Z$.  Over $\F_3$, both intermediate groups vanish.
These are exact coefficient changes of the same integer chain complex.

\subsection{Product torus fillings in every degree}

Fix $q\ge0$ and form
\[
 X_q=M\times\T^q
   =(V_0\times\T^q)\cup_{T\times\T^q}(V_1\times\T^q),
\]
using the identity on the auxiliary torus in the gluing.
This specifies a manifold of dimension $q+3$ and its complete boundary
identification.  When $q=3$, its regular Seifert fiber times $\T^3$
is a real four-torus.  The auxiliary circles are globally fixed.

Let $W=\Z^q$ with ordered basis $y_1,\ldots,y_q$.
Write $\Lambda^j W$ for the free abelian group on symbols
$y_{i_1}\wedge\cdots\wedge y_{i_j}$ with increasing indices.
Exchanging adjacent indices changes the sign, and repeated indices
give zero.  There are $\binom qj$ basis symbols, with
$\binom qj=0$ outside $0\le j\le q$.

The torus product cells have precisely these generators and zero
cellular boundary. Each circle edge has zero boundary.
For an oriented product,
\[
 \partial(a\times b)=(\partial a)\times b+(-1)^{\deg a}a\times\partial b.
\]
This formula follows by orienting each face with its outward normal first.

\begin{theorem}[All product homology groups]\label{prod:thm:product}
There is a cellular chain homotopy model
\begin{equation}\label{prod:eq:product-chain}
 C_*(X_q)=C_*(M)\otimes_{\Z}\Lambda^*W,
 \qquad \partial=\partial_M\otimes1,
\end{equation}
where $C_*(M)$ is \eqref{prod:eq:chain}, and total degrees add.
Consequently, if $\Delta\ne0$,
\begin{equation}\label{prod:eq:product-homology}
 H_j(X_q;\Z)\cong
 \Z^{\binom qj+\binom q{j-3}}
 \oplus(\Z/|\Delta|\Z)^{\binom q{j-1}}.
\end{equation}
If $\Delta=0$, then
\begin{equation}\label{prod:eq:product-zero}
 H_j(X_q;\Z)\cong
 \Z^{\binom qj+\binom q{j-1}+\binom q{j-2}+\binom q{j-3}}.
\end{equation}
The decomposition is induced by the fixed oriented product cells.
\end{theorem}

\begin{proof}
Take the product of the double mapping cylinder with $\T^q$.
Its cells are products of the cells used in Theorem~\ref{prod:thm:chains}
with the coordinate torus cells.  The displayed product boundary formula
gives the tensor differential.  Collapsing the contractible edge in the
first factor remains a homotopy equivalence after this product.
This proves \eqref{prod:eq:product-chain}.

The second factor has zero differential and a free basis in every degree.
Thus the total complex splits into copies of $C_*(M)$,
one shifted by $|I|$ for each subset $I\subseteq\{1,\ldots,q\}$.  Taking cycles and boundaries separately
in these direct summands proves
\[
 H_j(X_q;\Z)=\bigoplus_I H_{j-|I|}(M;\Z).
\]
Theorem~\ref{prod:thm:homology3} now gives both formulas.
No splitting of an unknown extension or omission of an integral torsion
term is needed.
\end{proof}

There is also a boundary description of every map in this calculation.
Write the boundary lattice as $\Z x\oplus\Z c\oplus W$ and the two
core lattices as $\Z z_i\oplus W$.  The inclusion maps are
\[
 J_0(x)=mz_0,\ J_0(c)=rz_0,\ J_0|_W=1,
 \qquad
 J_1(x)=-\nu z_1,\ J_1(c)=sz_1,\ J_1|_W=1.
\]
In degree $j$ the gluing map is
\begin{equation}\label{prod:eq:exterior}
 F_j=\Lambda^jJ_0\oplus(-\Lambda^jJ_1):
 \Lambda^j(\Z^2\oplus W)\longrightarrow
 \Lambda^j(\Z\oplus W)\oplus\Lambda^j(\Z\oplus W).
\end{equation}

On $\Lambda^jW$ it is the diagonal $(1,-1)$.
On $(\Z x\oplus\Z c)\wedge\Lambda^{j-1}W$ it is
$B\otimes1$.
It kills $(x\wedge c)\wedge\Lambda^{j-2}W$.
Indeed, both $x$ and $c$ map to multiples of the same core vector,
whose wedge with itself is zero.

These three formulas construct $F_2$ and $F_3$ explicitly, along
with every higher-degree map.  Canceling the diagonal summands in
the mapping-cylinder complex recovers \eqref{prod:eq:product-chain}.

\begin{example}[The six-dimensional case]\label{prod:ex:six}
For $q=3$ and $\Delta\ne0$, put $A=\Z/|\Delta|\Z$.
The complete answer is
\[
\begin{array}{c|ccccccc}
 j&0&1&2&3&4&5&6\\\hline
 H_j(X_3;\Z)&\Z&\Z^3\oplus A&\Z^3\oplus A^3&
 \Z^2\oplus A^3&\Z^3\oplus A&\Z^3&\Z.
\end{array}
\]
In Example~\ref{prod:ex:eight}, the three degree-two torsion generators
are $[e_1]\times[y_i]$, and the three degree-three torsion generators
are $[e_1]\times[y_i\wedge y_j]$ for $i<j$.
Their order-eight bounding chains are $e_2\times y_i$ and
$e_2\times(y_i\wedge y_j)$ in the reduced complex.
The two free degree-three classes are the oriented fundamental class
of $M$ times a point, and a point times the fundamental class of $\T^3$.

For effective data $(3,4;1,-1;0)$, one has $\Delta=1$.
The torsion vanishes, but these two degree-three classes and the three
degree-one classes remain.  Thus the unit determinant does not produce
the homology of a six-sphere in this explicit four-torus-fiber model.
\end{example}

\begin{proposition}[Naturality under torus maps]\label{prod:prop:naturality}
An integer matrix $L:\Z^q\to\Z^{q'}$ defines
$\ell:\T^q\to\T^{q'}$ by reduction modulo $\Z^{q'}$.
The product map $1_M\times\ell$ acts in the decomposition above as
$1\otimes\Lambda^jL$ on each torus degree $j$.
These maps compose with matrix multiplication and commute with all
integral boundaries and coefficient connecting maps.
\end{proposition}

\begin{proof}
On a coordinate circle the map has winding vector equal to the
corresponding column of $L$.
On an oriented coordinate $j$-torus, the coefficient of a target
$j$-cell is the determinant of the corresponding minor.
This determinant is the signed degree of the linear map between the two
coordinate tori.
Expanding this determinant is exactly the definition of $\Lambda^jL$.
The torus cellular boundaries vanish, so these homology coefficients
also specify its cellular chain map after cellular approximation.

Tensoring with the identity gives a chain map for
\eqref{prod:eq:product-chain}.  The determinant expansion proves the exterior
power composition law.  Lifting a coefficient cycle integrally and
dividing its boundary commutes with this integer chain map, proving the
last assertion.
\end{proof}

For $p\mid\Delta\ne0$, the mod-$p$ dimension in total degree $j$ is
\[
 \binom qj+\binom q{j-1}+\binom q{j-2}+\binom q{j-3}.
\]
The extra classes $[e_2]\otimes y_I$ map to
$(\Delta/p)[e_1]\otimes y_I$ under the integral connecting map.
For $p\nmid\Delta$, only the first and last binomial terms remain.
This coefficient dependence is determined on chains before taking homology.

\subsection{Nonfree local actions and the actual meridian}

The coprimality condition has a geometric meaning that survives beyond
formal presentations.  Let $r>0$, $m\in\Z$, and
$\zeta_r=\exp(2\pi i/r)$.  On $D^2\times S^1$ let the cyclic group
of order $r$ act by
\begin{equation}\label{prod:eq:local-action}
                        (z,w)\longmapsto(\zeta_rz,\zeta_r^m w).
\end{equation}
The action preserves the product orientation.  Its quotient maps to the
base disk by $t=z^r$.

Away from the core $z=0$ the action is free.
On the core its stabilizer has order $a=\gcd(r,m)$, where
$\gcd(r,0)=r$.

\begin{proposition}[Coarse quotient and primitive reduction]\label{prod:prop:local-reduction}
Put $R=r/a$ and $M=m/a$.  The underlying quotient of
\eqref{prod:eq:local-action} is an oriented solid torus with actual meridian
$Rx-Mc$ in its boundary marking $(x,c)$.
Its effective fiber multiplicity is $R$.
The original action is free exactly when $a=1$.
\end{proposition}

\begin{proof}
The stabilizer subgroup of order $a$ is generated by the $R$-th power
of the displayed generator.  It rotates $z$ by an $a$-th root of unity
and fixes $w$.  Its quotient is explicitly $(Z,w)=(z^a,w)$.

The residual group has order $R$ and acts by
$(Z,w)\mapsto(\zeta_RZ,\zeta_R^M w)$.
Since $\gcd(R,M)=1$, no nonidentity power fixes a core point.
Off the core, a fixed point already forces the disk rotation to be
trivial.  Hence the residual action is free.

Choose $h\in\Z$ with $hM\equiv1\pmod R$.
The map
\[
                  (Z,w)\longmapsto (Zw^{-h},w^R)
\]
is invariant under the residual action and takes values in $D^2\times S^1$.
It is onto.  Two preimages differ by that action. Their second coordinates differ by
an $R$-th root, and their first coordinates determine the same group element.

It therefore induces a homeomorphism from the
compact quotient to $D^2\times S^1$.
It preserves orientation, since the circle degree $R$ is positive and
the disk coordinate changes by rotations.

For the boundary marking, use angular coordinates for $z,w$ before
quotienting.  A lifted base loop has displacement $(1/r,m/r)$ and
represents $x$.  The regular fiber has displacement $(0,1)$ and
represents $c$.  After taking $Z=z^a$, the disk meridian has displacement
$(1/a,0)=R(1/r,m/r)-M(0,1)$.
Thus it represents $Rx-Mc$.
The base coordinate is $t=Z^R$, proving its multiplicity.
\end{proof}

For arbitrary data \eqref{prod:eq:data}, set
\begin{equation}\label{prod:eq:reduced-data}
 a_0=\gcd(r,m),\quad a_1=\gcd(s,n),\quad
 R=r/a_0,\ M=m/a_0,\ S=s/a_1,\ N=n/a_1.
\end{equation}
Gluing the two underlying quotient solid tori therefore uses
$\bar\mu_0=Rx-Mc$ and
$\bar\mu_1=Sx+(N-Sb_0)c$.
These are primitive, and their determinant is
\[
                 \bar\Delta=\Delta/(a_0a_1)\in\Z.
\]

The actual coarse manifold $\bar M$ consequently has
$H_1(\bar M)=\Z/|\bar\Delta|$ and $H_2(\bar M)=0$ when
$\Delta\ne0$.  At $\Delta=0$ it has $H_1=H_2=\Z$.
In both cases $H_3(\bar M)=\Z$.
These statements apply Theorem~\ref{prod:thm:homology3} to the primitive slopes.

\subsection{The unreduced gluing space and its comparison map}

The original integer data still define maps \eqref{prod:eq:core-maps} and
a double mapping cylinder $P_B$, even when a winding map is not onto
on fundamental groups.  Its ordinary fibers then have several circle
components.  The earlier cylinder argument no longer identifies each
half with a solid torus.  Nevertheless, its cellular complex remains
\eqref{prod:eq:chain}, since the prism calculation used only the specified
integer windings.

\begin{theorem}[Formal gluing and its coarse comparison]\label{prod:thm:nonfree}
Let $d=\gcd(r,s,m,n)$.
For the formal space $P_B$, one has $H_0(P_B)=H_3(P_B)=\Z$ and
\[
 (H_1(P_B),H_2(P_B))=
 \begin{cases}
 (\Z/d\Z\oplus\Z/(|\Delta|/d)\Z,0),&\Delta\ne0,\\
 (\Z\oplus\Z/d\Z,\Z),&\Delta=0.
 \end{cases}
\]
There is an explicit continuous map $\Phi:\bar M\to P_B$, after
using the cylinder realization of $\bar M$, with
\begin{equation}\label{prod:eq:comparison}
 \Phi_0=1,\quad \Phi_1=D,\quad\Phi_2=1_{\Z^2},\quad\Phi_3=1,
 \qquad D=\operatorname{diag}(a_0,a_1).
\end{equation}
It induces isomorphisms on $H_0,H_2,H_3$ and an exact sequence
\begin{equation}\label{prod:eq:nonfree-h1}
 0\longrightarrow H_1(\bar M)\xrightarrow{\Phi_*}H_1(P_B)
 \longrightarrow\Z/a_0\Z\oplus\Z/a_1\Z\longrightarrow0.
\end{equation}
For the relative homology of this map, the only nonzero group is
\[
                 H_1(\Phi;\Z)=\Z/a_0\Z\oplus\Z/a_1\Z.
\]
\end{theorem}

\begin{proof}
Integer row and column operations from the Euclidean algorithm put a
nonzero two-by-two matrix into Smith form
$\operatorname{diag}(d_1,d_2)$ with $d_1\mid d_2$.
Place a least positive nonzero entry at the upper left.
Divide each entry in its row and column by it with remainder.
A nonzero remainder gives a smaller pivot, so this process terminates.
If the pivot fails to divide the last diagonal entry, adding the second
column to the first produces another smaller remainder.
At termination the pivot divides all entries and equals their greatest
common divisor.

Clearing its row and column gives the diagonal form.
Unimodular operations preserve the entry gcd and the absolute determinant,
so $d_1=d$ and $d_1d_2=|\Delta|$ when $\Delta\ne0$.
For $\Delta=0$, the rank is one because $r,s>0$, and the second
entry is zero.  Applying this reduction to \eqref{prod:eq:chain} proves
the formal homology formulas.

Let $\bar f_i$ be the primitive winding maps with the reduced data.
Then $f_i=p_{a_i}\bar f_i$, where
$p_a:S^1\to S^1$ is multiplication by $a$.
Use the identity on $T\times[0,1]$ and these degree-$a_i$ maps on
the two cores.  The endpoint identifications commute exactly, so this
defines a continuous map between the two double mapping cylinders.

On cellular chains it is \eqref{prod:eq:comparison}, and
\[
                         B=D\bar B.
\]
Thus the map is injective in every chain degree.
Its quotient complex is $Q_1=\Z/a_0\oplus\Z/a_1$, zero in all
other degrees.

This proves that the degree-two kernels agree and the
degree-three and degree-zero groups agree.
In degree one, an element $Du$ is a boundary $Bv$ precisely when
$Du=D\bar Bv$, equivalently $u=\bar Bv$, proving injectivity.
The remaining quotient is $Q_1$, which proves
\eqref{prod:eq:nonfree-h1} directly.

The relative chain complex of a map $f:C\to C'$ is its cone,
$\operatorname{Cone}(f)_j=C'_j\oplus C_{j-1}$, with boundary
$(b,a)\mapsto(\partial b+f(a),-\partial a)$.
For an injective chain map, this cone maps onto the quotient complex by
$(b,a)\mapsto[b]$.  Its kernel identifies with pairs $(f(c),a)$
and has differential $(c,a)\mapsto(\partial c+a,-\partial a)$.
The homotopy $(c,a)\mapsto(0,c)$ contracts this kernel.
The cone and quotient therefore have the same homology.

The topological mapping cylinder of $\Phi$, relative to its source,
has this cone complex by the same prism calculation.
This proves the relative assertion.
\end{proof}

Tensoring \eqref{prod:eq:comparison} with $\Lambda^*\Z^q$ gives a
continuous product map $\Phi\times1_{\T^q}$.
If $A_0=\Z/a_0\oplus\Z/a_1$, its relative groups are
\begin{equation}\label{prod:eq:relative-product}
                   H_j(\Phi\times1;\Z)=A_0^{\binom q{j-1}}.
\end{equation}
There are short exact sequences
\[
 0\to H_j(\bar M\times\T^q)\to H_j(P_B\times\T^q)
          \to A_0^{\binom q{j-1}}\to0.
\]
They follow by taking the direct sum of the shifted degree-one sequence
\eqref{prod:eq:nonfree-h1}, together with the other homology isomorphisms.

\begin{example}\label{prod:ex:nonfree}
For $(r,s;m,n;b_0)=(2,4;2,2;0)$, one has
$a_0=a_1=2$ and $\Delta=12$.
The formal group is $\Z/2\oplus\Z/6$, while the coarse filling has
reduced data $(1,2;1,1;0)$ and first homology $\Z/3$.
The comparison quotient is $(\Z/2)^2$.
For the six-dimensional product, the formal degree-two group is
\[
 \Z^3\oplus(\Z/2)^3\oplus(\Z/6)^3,
\]
while the actual coarse degree-two group is $\Z^3\oplus(\Z/3)^3$.
The relative degree-two and degree-three groups are both $(\Z/2)^6$.

At $(2,4;2,-4;0)$, the determinant is zero and both reduced orders
are one.  The formal first homology is $\Z\oplus\Z/2$,
whereas the coarse first homology is $\Z$.
Here $a_0=2$, $a_1=4$, so the comparison quotient is
$\Z/2\oplus\Z/4$.

In core generators $z_0,z_1$, the formal relation is
$2z_0-4z_1=0$.  Put $t=z_0-2z_1$, so $2t=0$.
The coarse generator maps to $2z_0=4z_1$.
Thus the comparison multiplies the free coordinate by four and misses
the order-two summand.
Even at rank drop, effectivization changes an integral subgroup
and its torsion rather than merely changing notation.
\end{example}

\subsection{The information beyond the boundary matrix}

The additive complex retains higher homology, but it does not determine
the oriented linking of torsion cycles.  We give an explicit distinction
because it limits what one can conclude from an exact chain calculation.

Use the linking convention of Theorem~\ref{thm:linking},
$\lambda([z],[w])=(C\cdot w)/N\pmod{\Z}$ when $\partial C=Nz$.
Changing $C$ by a closed two-cycle does not change this value: its
integer intersection with a torsion class is zero.
Changing a cycle representative changes the numerator by a multiple
of $N$.  Thus the pairing is well defined.  Oriented homeomorphisms
preserve it because they preserve boundaries and intersection signs.

For $\mu_0=x$ and $\mu_1=ux+vc$, with $v\ne0$ and
$\gcd(u,v)=1$, the longitude $c$ generates $H_1=\Z/|v|$.
The meridian-chain computation in the proof of
Theorem~\ref{thm:linking} gives
\begin{equation}\label{prod:eq:linking-distinction}
                            \lambda(c,c)=-u/v\pmod\Z.
\end{equation}

Take $(u,v)=(1,5)$ and $(2,5)$.
Both fillings have the reduced additive complex
$\partial e_2=5e_1$, with all other boundaries zero.
Their linking values on the displayed generators are $-1/5$ and
$-2/5$.  Replacing a generator of $\Z/5$ multiplies its value by
a unit square.  The unit squares modulo five are one and four, so these
pairings are not isomorphic.  In particular the two oriented manifolds
are not orientation-preservingly homeomorphic, despite their identical
additive chain models.

A different limitation concerns the boundary maps themselves.
Theorem~\ref{prod:thm:product} assumes product gluing and globally fixed
auxiliary circles.  A torus fibration with nontrivial monodromy or an
additional degeneration has different neighborhoods and specialization
maps.  Its integral degree-two and degree-three calculation requires
those chain maps and their attachments.  Neither a two-generator
presentation nor a unit determinant specifies them.
The explicit product example with $\Delta=1$ has
$H_1=\Z^3$ and $H_3=\Z^2$, so it already separates these questions.
